\begin{textsample}{POS Dim 5 – pi – Score 9.00 – pi/pi\_ef\_27.txt}  \label{ex:f5_pos_231}
MATEMÁTICA NA BNCC Segundo a BNCC, o Ensino Fundamental deve ter compromisso com o desenvolvimento do \textbf{letramento} matemático, definido como as competências e habilidades de raciocinar, representar, comunicar e argumentar matematicamente de modo a favorecer o estabelecimento de conjecturas, a formulação e a resolução de problemas em uma variedade de contextos, utilizando conceitos, procedimentos, fatos e ferramentas matemáticas.
Ou seja, possui uma visão integral e o propósito de contribuir para a construção de uma sociedade mais ética, democrática, responsável, inclusiva, sustentável e solidária, que respeita e promove a diversidade e os direitos humanos, sem preconceitos de qualquer natureza.
COMPETÊNCIAS ESPECÍFICAS DA MATEMÁTICA ENQUANTO COMPONENTE CURRICULAR O componente de Matemática pretende firmar um compromisso com o desenvolvimento integral de todos os estudantes nas suas dimensões afetiva, ética, física, intelectual, moral, social e simbólica visando ser a mola propulsora para a melhoria da qualidade das aprendizagens de todos aqueles matriculados nas diferentes redes de ensino do Piauí.
É inegável a importância do componente, no entanto, vale ressaltar que a prática pedagógica deverá ser diferenciada a partir dos Anos Iniciais, contando com apoio e participação de toda a comunidade escolar e outros profissionais que possam colaborar no processo de ensino e aprendizagem, permitindo o desenvolvimento das habilidades e o alcance das competências.
Considerando estes pressupostos e em articulação com as competências gerais da Educação Básica, a área da Matemática e, por consequência, o componente curricular de Matemática, deve garantir o desenvolvimento das competências específicas elencadas abaixo : 1.
Reconhecer que a Matemática é uma ciência humana, fruto das necessidades e preocupações de diferentes culturas, em diferentes momentos históricos, e é uma ciência viva, que contribui para solucionar problemas científicos e tecnológicos e para alicerçar descobertas e construções, inclusive com impactos no mundo do trabalho. 2.
Desenvolver o raciocínio lógico, o espírito de investigação e a capacidade de produzir argumentos convincentes, recorrendo aos conhecimentos matemáticos para compreender e atuar no mundo. 3.
Compreender as relações entre conceitos e procedimentos dos diferentes campos da Matemática ( Aritmética, Álgebra, Geometria, Estatística e Probabilidade ) e de outras áreas do conhecimento, sentindo segurança quanto à própria capacidade de construir e aplicar conhecimentos matemáticos, desenvolvendo a autoestima e a perseverança na busca de soluções. 4.
Fazer observações sistemáticas de aspectos quantitativos e qualitativos presentes nas práticas sociais e culturais, de modo a investigar, organizar, representar e comunicar informações relevantes, para interpretá -las e avaliá -las crítica e eticamente, produzindo argumentos convincentes. 5.
Utilizar processos e ferramentas matemáticas, inclusive tecnologias digitais disponíveis, para modelar e resolver problemas cotidianos, sociais e de outras áreas de conhecimento, validando estratégias e resultados. 6.
Enfrentar situações-problema em múltiplos contextos, incluindo -se situações imaginadas, não diretamente relacionadas com o aspecto prático-utilitário, expressar suas respostas e sintetizar conclusões, utilizando diferentes registros e linguagens ( gráficos, tabelas, esquemas, além de texto escrito na língua materna e outras linguagens para descrever algoritmos, como fluxogramas, e dados ). 7.
Desenvolver e/ou discutir projetos que abordem, sobretudo, questões de urgência social, com base em princípios éticos, democráticos, sustentáveis e solidários, valorizando a diversidade de opiniões de indivíduos e de grupos sociais, sem preconceitos de qualquer natureza. 8.
Interagir com seus pares de forma cooperativa, trabalhando coletivamente no planejamento e desenvolvimento de pesquisas para responder a questionamentos e na busca de soluções para problemas, de modo a identificar aspectos consensuais ou não na discussão de uma determinada questão, respeitando o modo de pensar dos colegas e aprendendo com eles.
O Componente Curricular de Matemática propõe cinco unidades temáticas : números ; álgebra ; geometria ; grandezas e medidas ; e probabilidade e estatística, que orientam a formulação das habilidades a serem desenvolvidas ao longo do ensino fundamental.
A unidade temática números tem como finalidade desenvolver o pensamento numérico, que implica o conhecimento de maneiras de quantificar atributos de objetos e de julgar e interpretar argumentos baseados em quantidades.
No processo da construção da noção de número, os alunos precisam desenvolver, entre outras, as ideias de aproximação, proporcionalidade, equivalência e ordem, noções fundamentais da matemática.
Para essa construção, é importante propor, por meio de situações significativas, sucessivas ampliações dos campos numéricos.
No estudo desses campos numéricos, devem ser enfatizados registros, usos, significados e operações.
Lembrando que cada unidade temática pode receber ênfase diferente dependendo do ano de escolarização.
Nos Anos Iniciais do Ensino Fundamental em Números deverá acontecer o desenvolvimento de habilidades de : - Leitura, escrita e ordenação de números naturais e de números racionais por meio da identificação e compreensão de características do sistema de \textbf{numeração} decimal, sobretudo o valor posicional dos algarismos. - Elaboração e resolução de problemas com números naturais e números racionais cuja representação decimal é finita, envolvendo diferentes significados das operações. - Argumentação e justificativa dos procedimentos utilizados para a resolução e avaliação da plausibilidade dos resultados encontrados. - Cálculo por meio do desenvolvimento de diferentes estratégias para a obtenção dos resultados, sobretudo por estimativa e cálculo mental, além de algoritmos e uso de calculadoras.
Nos Anos Finais do Ensino Fundamental na Unidade temática Números, deverá acontecer o desenvolvimento de habilidades de : - Resolver problemas com números naturais, inteiros e racionais, envolvendo as operações fundamentais, com seus diferentes significados, e utilizando estratégias diversas, com compreensão dos processos neles envolvidos. - Aprofundar a noção de número, por meio da vivência de problemas, sobretudo os geométricos, nos quais os números racionais não são suficientes para resolvê -los, de modo que eles reconheçam a necessidade de outros números : os irracionais.
A unidade temática álgebra, por sua vez, tem como finalidade o desenvolvimento de um tipo especial de pensamento – pensamento algébrico – que é essencial para utilizar modelos matemáticos na compreensão, representação e análise de relações quantitativas de grandezas e, também, de situações e estruturas matemáticas, fazendo uso de \textbf{letras} e outros símbolos.
Para esse desenvolvimento, é necessário que os alunos identifiquem regularidades e padrões de \textbf{sequências} numéricas e não numéricas, estabeleçam leis matemáticas que expressem a relação de interdependência entre grandezas em diferentes contextos, bem como criar, interpretar e transitar entre as diversas representações gráficas e simbólicas, para resolver problemas por meio de equações e inequações, com compreensão dos procedimentos utilizados.
As ideias matemáticas fundamentais vinculadas a essa unidade são : equivalência, variação, interdependência e proporcionalidade.
Em síntese, essa unidade temática deve enfatizar o desenvolvimento de uma linguagem, o estabelecimento de generalizações, a análise da interdependência de grandezas e a resolução de problemas por meio de equações ou inequações.
Nos Anos Iniciais do Ensino Fundamental a relação dessa unidade temática com a de Números é bastante evidente no trabalho com : - \textbf{Sequências} – recursivas e repetitivas –, seja na ação de completar uma \textbf{sequência} com elementos ausentes, seja na construção de \textbf{sequências} segundo uma determinada regra de formação. - Cálculos do tipo 5 + \_\_\_\_\_ = 9. - A relação de equivalência pode ter seu início com atividades simples, envolvendo a igualdade, como reconhecer que se 2 + 3 = 5 e 5 = 4 + 1 então 2 + 3 = 4 + 1.
Atividades como essa contribuem para a compreensão de que o sinal de igualdade não é apenas a indicação de uma operação a ser feita. - A noção intuitiva de função pode ser explorada por meio da resolução de problemas envolvendo a variação proporcional direta entre duas grandezas ( sem utilizar a regra de três ), como : “ Se com duas medidas de suco concentrado eu obtenho três litros de refresco, quantas medidas desse suco concentrado eu preciso para ter doze litros de refresco? ” No Ensino Fundamental – Anos Finais, os estudos de Álgebra retomam, aprofundam e ampliam o que foi trabalhado no Ensino Fundamental – Anos Iniciais.
Nessa fase, os alunos devem compreender os diferentes significados das : variáveis numéricas em uma expressão, estabelecer uma generalização de uma propriedade, investigar a regularidade de uma \textbf{sequência} numérica, indicar um valor desconhecido em uma sentença algébrica e estabelecer a variação entre duas grandezas.
É necessário, portanto, que os alunos estabeleçam conexões entre variável e função e entre incógnita e equação.
As técnicas de resolução de equações e inequações, inclusive no plano cartesiano, devem ser desenvolvidas como uma maneira de representar e resolver determinados tipos de problema, e não como objetos de estudo em si mesmos.
A aprendizagem de Álgebra, como também aquelas relacionadas a outros campos da Matemática ( Números, Geometria e Probabilidade e Estatística ), podem contribuir para o desenvolvimento do pensamento computacional dos alunos, tendo em vista que eles precisam ser capazes de traduzir uma situação dada em outras linguagens, como transformar situações-problema, apresentadas em língua materna, em fórmulas, tabelas e gráficos e vice-versa.
A geometria envolve o estudo de um amplo conjunto de conceitos e procedimentos necessários para resolver problemas do mundo físico e de diferentes áreas do conhecimento.
Assim, nessa unidade temática, estudar posição e deslocamentos no espaço, formas e relações entre elementos de figuras planas e espaciais pode desenvolver o pensamento geométrico dos alunos.
Esse pensamento é necessário para investigar propriedades, fazer conjecturas e produzir argumentos geométricos convincentes.
Nos Anos Iniciais do Ensino Fundamental, a expectativa é de que os estudantes : - Identifiquem e estabeleçam pontos de referência para a localização e o deslocamento de objetos. - Construam representações de espaços conhecidos e estimem distâncias, usando, como suporte, mapas ( em papel, tablets ou smartphones ), croquis e outras representações.
Existem expectativas na unidade temática de geometria que são tanto dos Anos Iniciais como Finais, sendo necessário que os estudantes : - Indiquem características das formas geométricas tridimensionais e bidimensionais, associem figuras espaciais a suas planificações e vice-versa. - Nomeiem e comparem polígonos, por meio de propriedades relativas aos lados, vértices e ângulos. - Manipulem representações de figuras geométricas planas em quadriculados ou no plano cartesiano, e com recurso de softwares de geometria dinâmica para a identificação de simetrias. - Realizem, quando necessário, transformações entre unidades de medida padronizadas mais usuais. - Resolvam problemas sobre situações de compra e venda e desenvolvam, por exemplo, atitudes éticas e responsáveis em relação ao consumo.
Sugere -se que esse processo seja iniciado utilizando, preferencialmente, unidades de medidas não convencionais para fazer as comparações e \textbf{medições}.
Nos Anos Finais do Ensino Fundamental, a expectativa é que os estudantes : - Reconheçam comprimento, área, \textbf{volume} e abertura de ângulo como grandezas associadas a figuras geométricas. - Consigam resolver problemas envolvendo essas grandezas com o uso de unidades de medida padronizadas mais usuais. - Estabeleçam e utilizem relações entre essas grandezas e entre elas e grandezas não geométricas, para estudar grandezas derivadas como densidade, velocidade, energia, potência, entre outras.
Outro ponto a ser destacado refere -se à introdução de medidas de capacidade de armazenamento de computadores como grandeza associada a demandas da sociedade moderna.
Nesse caso, é importante destacar o fato de que os prefixos utilizados para byte ( quilo, mega, giga ) não estão associados ao sistema de \textbf{numeração} decimal, de base 10, pois um quilobyte, por exemplo, corresponde a 1.024 bytes, e não a 1.000 bytes.
As medidas quantificam grandezas do mundo físico e são fundamentais para a compreensão da realidade.
Assim, a unidade temática grandezas e medidas, ao propor o estudo das medidas e das relações entre elas, ou seja, das relações métricas, favorece a integração da matemática a outras áreas de conhecimento, como ciências ( densidade, grandezas e escalas do sistema solar, energia elétrica etc. ) ou geografia ( coordenadas geográficas, densidade demográfica, escalas de mapas e guias etc. ).
Essa unidade temática contribui ainda para a consolidação e ampliação da noção de número, a aplicação de noções geométricas e a construção do pensamento algébrico.
No Ensino Fundamental – Anos Iniciais, a expectativa é que os estudantes reconheçam que medir é comparar uma grandeza com uma unidade e expressar o resultado da comparação por meio de um número.
Além disso, devem resolver problemas oriundos de situações cotidianas que envolvam grandezas como comprimento, massa, tempo, temperatura, área ( de triângulos e retângulos ) e capacidade e \textbf{volume} ( de sólidos formados por \textbf{blocos} retangulares ), sem uso de fórmulas, recorrendo, quando necessário, a transformações entre unidades de medida padronizadas mais usuais.
No Ensino Fundamental – Anos Finais, a expectativa é a de que os estudantes reconheçam comprimento, área, \textbf{volume} e abertura de ângulo como grandezas associadas a figuras geométricas e que consigam resolver problemas envolvendo essas grandezas com o uso de unidades de medida padronizadas mais usuais.
Além disso, espera -se que estabeleçam e utilizem relações entre essas grandezas e entre elas e grandezas não geométricas, para estudar grandezas derivadas como densidade, velocidade, energia, potência, entre outras.
A incerteza e o tratamento de dados são estudados na unidade temática probabilidade e estatística.
Ela propõe a abordagem de conceitos, fatos e procedimentos presentes em muitas situações-problema da vida cotidiana, das ciências e da tecnologia, porque todos os cidadãos precisam desenvolver habilidades para coletar, organizar, representar, interpretar e analisar dados em uma variedade de contextos, de maneira a fazer julgamentos bem fundamentados e tomar as decisões adequadas.
Isso inclui raciocinar e utilizar conceitos, representações e índices estatísticos para descrever, explicar e predizer fenômenos.
No que concerne ao estudo de noções de probabilidade, busca -se promover a compreensão de que nem todos os fenômenos são determinísticos.
Para isso, o início da proposta de trabalho com probabilidade está centrado no desenvolvimento da noção de aleatoriedade, de modo que os alunos compreendam que há eventos certos, eventos impossíveis e eventos prováveis.
Merece destaque o uso de tecnologias – como calculadoras, para avaliar e comparar resultados, e planilhas eletrônicas, que ajudam na construção de gráficos e nos cálculos das medidas de tendência central.
A consulta a \textbf{páginas} de institutos de pesquisa – como a do Instituto Brasileiro de Geografia e Estatística ( IBGE ) – pode oferecer contextos potencialmente ricos não apenas para aprender conceitos e procedimentos estatísticos, mas também para utilizá -los com o intuito de compreender a realidade.
Se tratando de probabilidade, a finalidade, no Ensino Fundamental – Anos Iniciais, é promover a compreensão de que nem todos os fenômenos são determinísticos.
Para isso, o início da proposta de trabalho com probabilidade está centrado no desenvolvimento da noção de aleatoriedade, de modo que os alunos compreendam que há eventos certos, eventos impossíveis e eventos prováveis.
No Ensino Fundamental – Anos Finais, o estudo deve ser ampliado e aprofundado por meio de atividades nas quais os alunos façam experimentos aleatórios e simulações para confrontar os resultados obtidos com a probabilidade teórica – probabilidade frequentista.
A \textbf{progressão} dos conhecimentos se faz pelo aprimoramento da capacidade de enumeração dos elementos do espaço amostral, que está associada, também, aos problemas de contagem.
Com relação à estatística, os primeiros passos envolvem o trabalho com a coleta e a organização de dados de uma pesquisa de interesse dos alunos.
O planejamento de como fazer a pesquisa ajuda a compreender o papel da estatística no cotidiano dos alunos.
Assim, a leitura, a interpretação e a construção de tabelas e gráficos têm papel fundamental, bem como a forma de produção de texto escrito para a comunicação de dados, pois é preciso compreender que o texto deve sintetizar ou justificar as conclusões.
No Ensino Fundamental – Anos Finais, a expectativa é que os alunos saibam planejar e construir relatórios de pesquisas estatísticas descritivas, incluindo medidas de tendência central e construção de tabelas e diversos tipos de gráfico.
Esse planejamento inclui a definição de questões relevantes e da população a ser pesquisada, a decisão sobre a necessidade ou não de usar amostra e, quando for o caso, a seleção de seus elementos por meio de uma adequada técnica de amostragem.
Podemos definir que as unidades temáticas estão organizadas em agrupamentos, porém uma depende da outra para ser contemplada dentro dos objetos do conhecimento e das habilidades necessárias para a aprendizagem e não somente isso, mas uma habilidade depende de habilidades de outros componentes para serem contempladas, portanto a interdisciplinaridade é a forma mais eficaz de se efetivar uma verdadeira aprendizagem, utilizando uma diversidade metodológica e tecnológica de acordo com o público e com a realidade local da instituição.
A aprendizagem em matemática está intrinsecamente relacionada à compreensão, ou seja, à apreensão de significados dos objetos matemáticos, sem deixar de lado suas aplicações.
Em todas as unidades temáticas, a delimitação dos objetos de conhecimento e das habilidades considera que as noções matemáticas sejam retomadas, ampliadas e aprofundadas ano a ano.
Portanto, nos leva a compreender que, independentemente de o objeto do conhecimento aparecer ou não no ano seguinte, ele deve ser retomado para poder ampliar e aprofundar as habilidades, levando sempre o estudante a uma perspectiva de formulação de situação problemas em outros contextos.

% matched lemmas: bloco, letra, letramento, medição, numeração, progressão, página, sequência, volume
\end{textsample}
